\documentclass[11pt, letterpaper]{article}

% --- Idioma y codificación ---
\usepackage[spanish]{babel}
\usepackage[utf8]{inputenc}
\usepackage[T1]{fontenc}
\usepackage{lmodern}

% --- Matemáticas ---
\usepackage{amsmath, amssymb, amsthm}

% --- Formato ---
\usepackage[margin=2.5cm]{geometry}
\usepackage{parskip}
\usepackage{booktabs}
\usepackage{hyperref}
\hypersetup{colorlinks=true, linkcolor=blue!70!black, urlcolor=blue!70!black}

% --- Entornos de teoremas ---
\theoremstyle{definition}
\newtheorem{definition}{Definición}[section]
\theoremstyle{plain}
\newtheorem{theorem}{Teorema}[section]
\newtheorem{lemma}[theorem]{Lema}
\newtheorem{corollary}[theorem]{Corolario}
\theoremstyle{remark}
\newtheorem{remark}{Observación}[section]

\title{%
  T03 --- Demostración formal:\\[4pt]
  Prim produce un MST por la propiedad de corte
}
\author{Curso Linux--C--Rust --- B15 Estructuras de datos en C}
\date{}

\begin{document}
\maketitle
\tableofcontents

% ===================================================================
\section{Modelo}
% ===================================================================

\begin{definition}
Grafo $G = (V, E)$ no dirigido, conexo, con función de peso
$w: E \to \mathbb{R}$. $|V| = n$, $|E| = m$.
\end{definition}

\begin{definition}
\textbf{Algoritmo de Prim.} Mantiene un conjunto $S$ de vértices ya
incluidos en el MST parcial (inicialmente $S = \{s\}$ para un
vértice fuente arbitrario $s$). En cada paso, agrega la arista de
peso mínimo que cruza el corte $(S,\, V \setminus S)$, incorporando
su extremo $v \in V \setminus S$ a $S$.
\end{definition}

\textbf{Convención.} Pesos distintos (la prueba se adapta con
desempate consistente para pesos iguales).

% ===================================================================
\section{Demostración 1: correctitud por invariante de bucle}
% ===================================================================

\begin{theorem}\label{thm:correctness}
Al terminar, las aristas seleccionadas por Prim forman un MST.
\end{theorem}

\begin{proof}
Sea $A_k$ el conjunto de aristas seleccionadas tras $k$ iteraciones.
Definimos el invariante:
\[
  \mathrm{INV}(k):\quad
  A_k \text{ es subconjunto de algún MST de } G
\]

\textbf{Inicialización ($k = 0$).} $A_0 = \emptyset$, que es
subconjunto de cualquier MST. \checkmark

\textbf{Mantenimiento.} Supongamos $\mathrm{INV}(k)$: existe un MST
$T^*$ con $A_k \subseteq T^*$. En la iteración $k + 1$, Prim
selecciona la arista $e = (u, v)$ de peso mínimo que cruza el corte
$(S_k,\, V \setminus S_k)$, donde $S_k$ es el conjunto de vértices
incluidos tras $k$ iteraciones.

\emph{Caso 1: $e \in T^*$.} Entonces
$A_{k+1} = A_k \cup \{e\} \subseteq T^*$. \checkmark

\emph{Caso 2: $e \notin T^*$.} Agregar $e$ a $T^*$ crea un ciclo
$C$. El ciclo $C$ contiene $e = (u, v)$ con $u \in S_k$ y
$v \in V \setminus S_k$, así que $C$ cruza el corte al menos dos
veces. Sea $e' \neq e$ otra arista de $C$ que cruza el corte
$(S_k,\, V \setminus S_k)$, con $e' \in T^*$.

Como $e$ es la arista de peso mínimo que cruza el corte:
$w(e) < w(e')$ (pesos distintos).

Definimos $T^{**} = T^* \setminus \{e'\} \cup \{e\}$.

\begin{itemize}
  \item $T^{**}$ es un árbol generador (eliminar $e'$ del ciclo $C$
    y mantener $e$ preserva la conexión; $n - 1$ aristas).
  \item $w(T^{**}) = w(T^*) - w(e') + w(e) < w(T^*)$.
\end{itemize}

Pero $T^*$ era un MST, así que $w(T^{**}) \geq w(T^*)$.
Contradicción. $\therefore$ el Caso~2 no ocurre (con pesos
distintos).

Con pesos iguales: $w(e) \leq w(e')$, y $T^{**}$ es también un MST
con $A_{k+1} \subseteq T^{**}$. \checkmark

\textbf{Verificación de que $e' \notin A_k$.} La arista $e'$ cruza
el corte $(S_k,\, V \setminus S_k)$ y está en $T^*$. Si
$e' \in A_k$, entonces $e'$ fue seleccionada en una iteración
anterior $j < k$, cuando cruzaba el corte
$(S_j,\, V \setminus S_j)$. Pero entonces ambos extremos de $e'$
estarían en $S_k$ (el extremo de $V \setminus S_j$ fue incorporado a
$S$ al seleccionar $e'$). Esto contradice que $e'$ cruza
$(S_k,\, V \setminus S_k)$. $\therefore e' \notin A_k$, y el
intercambio es válido. \checkmark

\textbf{Terminación ($k = n - 1$).} $A_{n-1}$ contiene $n - 1$
aristas, $S = V$, y $A_{n-1}$ es subconjunto de algún MST. Como un
MST tiene exactamente $n - 1$ aristas, $A_{n-1}$ \textbf{es} un MST.
\end{proof}

% ===================================================================
\section{Demostración 2: key[v] es la arista ligera del corte}
% ===================================================================

\begin{theorem}\label{thm:key}
Al momento de extraer $v$ de la cola de prioridad,
$\mathrm{key}[v]$ es el peso de la arista de peso mínimo entre $v$
y algún vértice de $S$.
\end{theorem}

\begin{proof}
Se mantiene el invariante:
\[
  \mathrm{key}[v]
  = \min\{w(u, v) : u \in S,\; (u, v) \in E\}
\]
(o $\infty$ si no existe tal arista).

\emph{Inicialización.} $S = \{s\}$. Para cada vecino $v$ de $s$:
$\mathrm{key}[v] = w(s, v)$. Para no vecinos:
$\mathrm{key}[v] = \infty$. Correcto. \checkmark

\emph{Mantenimiento.} Al agregar $u$ a $S$, para cada vecino
$v \notin S$ de $u$: si $w(u, v) < \mathrm{key}[v]$, se actualiza
$\mathrm{key}[v] = w(u, v)$. Esto incorpora la nueva información de
que $u$ ahora está en $S$.

Ninguna arista más ligera entre $v$ y $S$ puede quedar sin
considerar: la única información nueva es que $u \in S$, y se examina
$w(u, v)$. Las aristas desde vértices anteriores de $S$ ya fueron
consideradas cuando esos vértices se agregaron. \checkmark
\end{proof}

\begin{corollary}
La arista seleccionada $(u, v)$ con
$\mathrm{key}[v] = \min_{v \notin S} \mathrm{key}[v]$ es la arista
de peso mínimo que cruza $(S,\, V \setminus S)$. Esto es exactamente
lo que requiere la propiedad de corte.
\end{corollary}

% ===================================================================
\section{Demostración 3: Prim y Dijkstra son estructuralmente
  idénticos}
% ===================================================================

\begin{theorem}\label{thm:equiv}
Prim y Dijkstra tienen la misma estructura algorítmica; difieren
únicamente en la clave de la cola de prioridad.
\end{theorem}

\begin{proof}
Comparación directa:

\begin{table}[h]
\centering
\begin{tabular}{lll}
\toprule
\textbf{Aspecto} & \textbf{Dijkstra} & \textbf{Prim} \\
\midrule
Cola de prioridad
  & $\mathrm{dist}[v] = \delta(s,v)$ estimada
  & $\mathrm{key}[v] = \min w(u,v)$, $u \in S$ \\
Actualización
  & $\min(\mathrm{dist}[v],\, \mathrm{dist}[u]{+}w(u,v))$
  & $\min(\mathrm{key}[v],\, w(u,v))$ \\
Optimiza
  & Distancia acumulada desde $s$
  & Peso de una sola arista \\
Resultado
  & Árbol de caminos más cortos
  & Árbol generador mínimo \\
\bottomrule
\end{tabular}
\end{table}

La diferencia es una sola línea: Dijkstra acumula
($\mathrm{dist}[u] + w$); Prim no ($w$ solo). Ambos extraen el
mínimo de la cola y relajan vecinos. La estructura de bucle, las
operaciones sobre el heap y el patrón de visita son idénticos.
\end{proof}

\begin{corollary}
La complejidad de Prim es idéntica a la de Dijkstra: $O(n^2)$ con
array, $O(m \log n)$ con heap binario, $O(m + n\log n)$ con
Fibonacci heap.
\end{corollary}

% ===================================================================
\section{Demostración 4: complejidad
  \texorpdfstring{$O(m \log n)$}{O(m log n)} con heap binario}
% ===================================================================

\begin{theorem}\label{thm:complexity}
Prim con min-heap binario (y lazy deletion) ejecuta en
$O(m \log n)$.
\end{theorem}

\begin{proof}
Contamos operaciones sobre el heap:

\textbf{Extracciones.} Se extraen a lo sumo $n$ vértices válidos y
hasta $m - n$ entradas obsoletas (lazy deletion). Total: $O(m)$
extracciones, cada una $O(\log m) = O(\log n)$:
\[
  C_{\text{extract}} = O(m \log n)
\]

\textbf{Inserciones.} Cada relajación exitosa
($w(u,v) < \mathrm{key}[v]$) produce un push al heap. Se procesan
$O(m)$ pares $(u, v)$ en total. Inserciones: $O(m)$, cada una
$O(\log n)$:
\[
  C_{\text{insert}} = O(m \log n)
\]

\textbf{Total:}
\[
  T = O(m \log n)
\]

\textbf{Con decrease-key (heap indexado):} a lo sumo $m$
decrease-keys, cada una $O(\log n)$, y $n$ extracciones
$O(\log n)$:
\[
  T = O((m + n) \log n) = O(m \log n)
\]
Mismo orden asintótico, menor constante por no tener entradas
duplicadas.
\end{proof}

% ===================================================================
\section{Demostración 5: Prim con array es óptimo para grafos densos}
% ===================================================================

\begin{theorem}\label{thm:dense}
Para grafos con $m = \Theta(n^2)$, Prim con array lineal ($O(n^2)$)
es asintóticamente óptimo.
\end{theorem}

\begin{proof}
\emph{Cota superior.} Ya demostrada: $O(n^2)$ con array.

\emph{Cota inferior.} Cualquier algoritmo MST debe examinar todas
las aristas (de lo contrario, una arista no examinada de peso
$-\infty$ invalidaría el resultado). Para grafos densos con
$m = \Theta(n^2)$:
\[
  T \geq \Omega(m) = \Omega(n^2)
\]

Combinando: $T = \Theta(n^2)$ para Prim con array en grafos densos.

\textbf{Comparación con heap para grafos densos ($m = n^2$):}
\[
  T_{\text{array}} = O(n^2), \qquad
  T_{\text{heap}} = O(n^2 \log n)
\]

El array gana por un factor de $\log n$. Para $n = 10^4$ con grafo
completo: $n^2 = 10^8$ vs $n^2 \log n \approx 1.3 \times 10^9$
--- el heap es ${\sim}\,13\times$ más lento.
\end{proof}

% ===================================================================
\section{Verificación numérica}
% ===================================================================

Grafo con 6 vértices, fuente~0:

\begin{table}[h]
\centering
\begin{tabular}{cccccc}
\toprule
\textbf{Iter.} & \textbf{Extraer $v$} & $\mathrm{key}[v]$
  & \textbf{Arista MST} & \textbf{Corte $(S, V \setminus S)$}
  & \textbf{¿Mín.?} \\
\midrule
1 & 0 & 0 & --- & $(\{0\},\, \{1,2,3,4,5\})$ & --- \\
2 & 3 & 3 & $(0,3)$
  & $(\{0\},\, \{1,2,3,4,5\})$ & Sí: $3 < 4$ \checkmark \\
3 & 1 & 4 & $(0,1)$
  & $(\{0,3\},\, \{1,2,4,5\})$ & Sí: $4 < 8,9$ \checkmark \\
4 & 2 & 2 & $(1,2)$
  & $(\{0,1,3\},\, \{2,4,5\})$ & Sí: $2 < 5,6,7,9$ \checkmark \\
5 & 4 & 5 & $(2,4)$
  & $(\{0,1,2,3\},\, \{4,5\})$ & Sí: $5 < 6,7,9$ \checkmark \\
6 & 5 & 1 & $(4,5)$
  & $(\{0,1,2,3,4\},\, \{5\})$ & Sí: $1 < 7$ \checkmark \\
\bottomrule
\end{tabular}
\end{table}

Peso total: $3 + 4 + 2 + 5 + 1 = 15$. Idéntico a Kruskal.
\checkmark

En cada paso, la arista seleccionada es efectivamente la de peso
mínimo que cruza el corte, confirmando la propiedad de corte.
\checkmark

% ===================================================================
\section{Resumen de resultados}
% ===================================================================

\begin{table}[h]
\centering
\begin{tabular}{llp{5.5cm}}
\toprule
\textbf{Resultado} & \textbf{Técnica} & \textbf{Clave} \\
\midrule
Prim produce MST
  & Invariante $A_k \subseteq T^*$ + corte
  & Intercambio $T^{**} = T^* \setminus \{e'\} \cup \{e\}$ \\
$\mathrm{key}[v]$ = arista ligera
  & Invariante de la cola
  & Actualizar al agregar cada $u$ a $S$ \\
Prim $\equiv$ Dijkstra
  & Comparación directa
  & Diferencia: $w(u,v)$ vs $\mathrm{dist}[u] + w(u,v)$ \\
$O(m \log n)$ con heap
  & Contar operaciones
  & $O(m)$ ins./ext.\ $\times\, O(\log n)$ \\
$\Theta(n^2)$ óptimo densos
  & Cota inferior $\Omega(m)$
  & Toda arista debe examinarse \\
\bottomrule
\end{tabular}
\end{table}

\end{document}
